%% LyX 2.4.3 created this file.  For more info, see https://www.lyx.org/.
%% Do not edit unless you really know what you are doing.
\documentclass[11pt]{article}
\renewcommand{\familydefault}{\sfdefault}
\usepackage{textcomp}
\usepackage[utf8]{inputenc}
\usepackage{float}
\usepackage{wrapfig}
\usepackage{booktabs}
\usepackage{varwidth}
\usepackage{amsmath}
\usepackage{amssymb}
\usepackage{graphicx}
\usepackage[a4paper]{geometry}
\geometry{verbose,tmargin=2cm,bmargin=2cm,lmargin=2cm,rmargin=2cm}
\usepackage{rotfloat}

\makeatletter

%%%%%%%%%%%%%%%%%%%%%%%%%%%%%% LyX specific LaTeX commands.
\newcommand*\LyXZeroWidthSpace{\hspace{0pt}}
%% Because html converters don't know tabularnewline
\providecommand{\tabularnewline}{\\}
%% Variable width box for table cells
\newenvironment{cellvarwidth}[1][t]
    {\begin{varwidth}[#1]{\linewidth}}
    {\@finalstrut\@arstrutbox\end{varwidth}}

\@ifundefined{date}{}{\date{}}
%%%%%%%%%%%%%%%%%%%%%%%%%%%%%% User specified LaTeX commands.
\usepackage{graphicx}
\usepackage[table]{xcolor}
\usepackage{caption}

\captionsetup{
    font=small,          % taille du texte de la légende
    labelfont=bf         % "Figure/Table X" en gras
}

\makeatother

\begin{document}
\title{{\bfseries{}%
\begin{minipage}[c]{5cm}%
\textbf{\includegraphics[width=0.9\columnwidth]{logo_lpcno_600_dpi-C}}%
\end{minipage}}{\bfseries{}%
\begin{minipage}[c]{2cm}%
\textbf{\includegraphics[width=0.9\columnwidth]{/home/romuald/ANALYSE/pyPCPrivate/gSOMOS/somos/config/svg/SOMOs_logo_vectorized}}%
\end{minipage}}{\bfseries{}%
\begin{minipage}[c]{0.5\columnwidth}%
\begin{center}
\textbf{SOMOs Identification}\\
\textbf{\rule[0.5ex]{1\columnwidth}{1pt}}
\par\end{center}%
\end{minipage}}}
\author{{\footnotesize Romuald Poteau, romuald.poteau@univ-tlse3.fr, LPCNO}}
\date{{\footnotesize\today}}
\maketitle

\section{Introduction}

To identify singly occupied molecular orbitals (\textbf{SOMOs}) in
open-shell systems, we employed two complementary schemes: (i) orbital
projection analysis, and (ii) cosine similarity mapping. Both approaches
compare the sets of $\alpha$ and $\beta$ molecular orbitals obtained
from unrestricted calculations. The projection scheme evaluates the
squared overlap of each $\alpha$ orbital with the $\beta$ orbital
space using the atomic orbital overlap matrix, allowing for a quantitative
decomposition of each $\alpha$ orbital across the $\beta$ manifold.
In contrast, the cosine similarity approach measures the angular similarity
between $\alpha$ and $\beta$ orbitals based on their MO coefficients,
identifying pairs of orbitals with nearly identical spatial character.
Together, these methods help identify SOMOs as occupied $\alpha$
orbitals that lack a clear counterpart among the occupied $\beta$
orbitals. Instead, they show strong projections onto the $\beta$
virtual space and minimal overlap with the $\beta$ occupied space
--- a characteristic signature of magnetic (unpaired) molecular orbitals
in open-shell systems.

\section{Similarity schemes}

\subsection{Projection of occupied $\alpha$ MOs onto the $\beta$ orbital space}

Given a Gaussian log file from an unrestricted DFT calculation, we
extract the molecular orbital (MO) coefficients for both $\alpha$
and $\beta$ orbitals (\texttt{pop=full} keyword), along with the
AO overlap matrix $\boldsymbol{S}$ (\texttt{iop(3/33=1)} keyword)
and the printing of the basis set (\texttt{gfprint gfinput} keywords).
The analysis focuses on evaluating how each occupied $\alpha$ orbital
projects onto the full space spanned by all $\beta$ orbitals, which
includes both occupied and virtual ones.

Let $\boldsymbol{\Phi}_{i}^{\alpha}\in\mathbb{R}^{1\times n_{\text{basis}}}$
be the coefficient vector of the $i$-th occupied $\alpha$ orbital,
and let $\boldsymbol{\Phi}^{\beta}\in\mathbb{R}^{N\times n_{\text{basis}}}$
be the matrix of all $\beta$ orbitals stored row-wise, where $N=n_{\beta}$
is the total number of $\beta$ orbitals. The projection vector is
computed as: 
\[
\mathbf{v}_{i}=<\phi_{i}^{\alpha}|\phi^{\beta}>=\boldsymbol{\Phi}_{i}^{\alpha}\cdot\boldsymbol{S}\cdot(\boldsymbol{\Phi}^{\beta})^{T}\in\mathbb{R}^{1\times N}
\]
The squared norm $\|\mathbf{v}_{i}\|^{2}$ gives the total overlap
of the $\alpha$ orbital with the $\beta$ space.

To differentiate between the contributions from occupied and virtual
$\beta$ orbitals, we split the projection: 
\begin{align*}
\mathbf{v}_{i}^{\text{occ}} & =<\phi_{i}^{\alpha}|\phi_{\mathrm{occ}}^{\beta}>=\boldsymbol{\Phi}_{i}^{\alpha}\cdot\boldsymbol{S}\cdot(\boldsymbol{\Phi}_{\text{occ}}^{\beta})^{T}\\
\mathbf{v}_{i}^{\text{virt}} & =<\phi_{i}^{\alpha}|\phi_{\mathrm{virt}}^{\beta}>=\boldsymbol{\Phi}_{i}^{\alpha}\cdot\boldsymbol{S}\cdot(\boldsymbol{\Phi}_{\text{virt}}^{\beta})^{T}
\end{align*}
We then compute:

$\|\mathbf{v}_{i}^{\text{occ}}\|^{2}$ = projection of $|\phi_{i}^{\alpha}>$
onto occupied $\beta$ orbitals

$\|\mathbf{v}_{i}^{\text{virt}}\|^{2}$ = projection of $|\phi_{i}^{\alpha}>$
onto virtual $\beta$ orbitals

The total projection norm is decomposed to analyze how concentrated
or spread the projection is across $\beta$ orbitals: 
\begin{itemize}
\item The three largest values among the squared projections $v_{ij}^{2}$
are summed to compute \textbf{``Top 1 (\%)}'', \textbf{``Top 2
(\%)}'' and \textbf{``Top 3 (\%)}''. \textbf{Top 1 (\%)} can also
bee seen as a \textbf{dominance ratio}, \emph{i.e}. the quantity is
defined as the largest single squared projection divided by the total
projection norm: $\max_{j}v_{ij}^{2}/\|\mathbf{v}_{i}\|^{2}$
\item The ``\textbf{$\beta$ MOs \textgreater 15\%}'' column lists all
$\beta$ orbitals contributing more than the specified percentage
to the squared projection norm, along with their contribution in the
format $[j,p_{j}]$, where $j$ is the index (1-based) and $p_{j}$
the percentage contribution. For the most important contribution,
it is nothing else than the \textbf{dominance ratio}. It provides
a direct quantitative decomposition of each $\alpha$ orbital onto
the $\beta$ orbital basis. Each entry explicitly identifies the $\beta$
orbital(s) that significantly compose the corresponding $\alpha$
orbital, along with their respective percentage contributions
\end{itemize}
An orbital is flagged as a \textbf{SOMO candidate} if its projection
onto the virtual $\beta$ space exceeds 0.5 and its projection onto
the occupied $\beta$ space is below 0.5: 
\[
\|\mathbf{v}_{i}^{\text{virt}}\|^{2}>0.5\quad\text{and}\quad\|\mathbf{v}_{i}^{\text{occ}}\|^{2}<0.5
\]

This criterion is named \textbf{``SOMO P2v?}'' in the output and
in a saved spreadsheet file. In some cases, where the mixing of the
projection onto occupied and virtual $\beta$ MOs makes the identification
not straightforward, a secondary, less robust, criterion has been
defined. A SOMO candidate, named ``\textbf{SOMO dom. $\boldsymbol{\beta}$
MO?}'', is identified when the dominance ratio is associated to a
virtual $\beta$ MO.

\subsection{Diagonalization of the projection of $\alpha$ occupied orbitals
onto $\beta$ subspaces}

In order to further analyze the nature of singly occupied molecular
orbitals (SOMOs) and their relation to the $\beta$ spin manifold,
a complementary diagonalization procedure was implemented.

Starting from the set of occupied $\alpha$ orbitals $\phi_{i}^{\alpha}$,
two separate projections are constructed, namely the projection onto
the occupied $\beta$ orbitals $\phi_{\text{occ}}^{\beta}$ and the
projection onto the virtual $\beta$ orbitals $\phi_{\text{virt}}^{\beta}$. 

Given the atomic orbital overlap matrix $\boldsymbol{S}$, the rectangular
projection matrices are defined as: 
\[
\boldsymbol{A}_{\text{occ}}=\boldsymbol{\Phi}_{i}^{\alpha}\cdot\boldsymbol{S}\cdot(\boldsymbol{\Phi}_{\text{occ}}^{\beta})^{T},\quad\boldsymbol{A}_{\text{virt}}=\boldsymbol{\Phi}_{i}^{\alpha}\cdot\boldsymbol{S}\cdot(\boldsymbol{\Phi}_{\text{virt}}^{\beta})^{T}.
\]

From these, the symmetric projection matrices are formed: 
\[
\boldsymbol{P}_{\text{occ}}=\boldsymbol{A}_{\text{occ}}\boldsymbol{A}_{\text{occ}}^{T},\quad\boldsymbol{P}_{\text{virt}}=\boldsymbol{A}_{\text{virt}}\boldsymbol{A}_{\text{virt}}^{T}.
\]

The matrices $\boldsymbol{P}_{\text{occ}}$ and $\boldsymbol{P}_{\text{virt}}$
are diagonalized to obtain their eigenvalues and eigenvectors.

The eigenvalues of $\boldsymbol{P}_{\text{occ}}$ quantify how strongly
a linear combination of occupied $\alpha$ orbitals projects onto
the occupied $\beta$ space. Similarly, the eigenvalues of $\boldsymbol{P}_{\text{virt}}$
measure the projection onto the virtual $\beta$ space. Eigenvectors
with low eigenvalues for $\boldsymbol{P}_{\text{occ}}$ but significant
projection onto $\beta$ virtual orbitals are strong candidates for
SOMOs.

\subsection{Cosine similarity of MOs}

The identification of singly occupied molecular orbitals (SOMOs) can
also be achieved through the computation of the cosine similarity
between pairs of molecular orbitals (MOs) derived from unrestricted
spin density functional theory (DFT) calculations. Specifically, we
computed similarities between $\alpha$ and $\beta$ spin orbitals,
taking into account the non-orthogonality of the basis set used in
quantum chemical calculations.

Let us denote two molecular orbital coefficient vectors as $\boldsymbol{\Phi}_{i}^{\alpha}$
(for alpha-spin orbitals) and $\boldsymbol{\Phi}_{j}^{\beta}$ (for
beta-spin orbitals). Each vector has dimensions corresponding to the
number of basis functions used in the calculation, denoted by $n_{\text{basis}}$.
Given the overlap matrix $\boldsymbol{S}$ (dimension $n_{\text{basis}}\times n_{\text{basis}}$),
obtained from the quantum chemistry calculation, the scalar product
between two coefficient vectors accounting for basis overlap is defined
as:

\begin{equation}
\langle\phi_{i}^{\alpha}|\phi_{j}^{\beta}\rangle=\boldsymbol{\Phi}_{i}^{\alpha T}\boldsymbol{S}\boldsymbol{\Phi}_{j}^{\beta}
\end{equation}

Thus, the cosine similarity between two molecular orbitals $\boldsymbol{\Phi}_{i}^{\alpha}$
and $\boldsymbol{\Phi}_{j}^{\beta}$ accounting for the basis overlap
matrix, $\boldsymbol{S}$, is given by:

\begin{equation}
\text{cosine similarity}(\phi_{i}^{\alpha},\phi_{j}^{\beta})=\frac{\boldsymbol{\Phi}_{i}^{\alpha T}\boldsymbol{S}\Phi_{j}^{\beta}}{\sqrt{\boldsymbol{\Phi}_{i}^{\alpha T}\boldsymbol{S}\boldsymbol{\Phi}_{i}^{\alpha}}\sqrt{\boldsymbol{\Phi}_{j}^{\beta T}\boldsymbol{S}\boldsymbol{\Phi}_{j}^{\beta}}}
\end{equation}

The similarity matrix constructed from these cosine similarities was
then used to optimally match $\alpha$ and $\beta$ orbitals employing
the Hungarian algorithm, ensuring maximal global similarity. Also
known as the Kuhn--Munkres algorithm, it is a classic method used
to solve the assignment problem: given a cost matrix, it finds the
optimal one-to-one assignment (or matching) that minimizes (or maximizes)
the total cost (or similarity). It is provided by the \texttt{scipy.optimize.linear\_sum\_assignment}
function. 

Orbital pairs with high cosine similarity, particularly those involving
occupied alpha-spin orbitals matched to virtual beta-spin orbitals
(or vice versa), can also be identified as potential candidates for
SOMOs. This method provides a robust and quantitatively precise approach
to identifying SOMOs in unrestricted DFT calculations, facilitating
detailed analyses of electronic structures in open-shell systems.

\subsection{Brief discussion}

The projection technique quantifies how much each $\alpha$ orbital
overlaps with the entire $\beta$ orbital space by computing the squared
norm of the projection vector using the AO overlap matrix. This provides
an absolute, physically meaningful measure of orbital mixing, especially
relevant when analyzing partial spin contamination or magnetic character.
In contrast, cosine similarity evaluates the angle between two orbital
vectors, yielding a dimensionless similarity score between -1 and
1. It's more suited for comparing the shape of orbitals than their
actual physical contribution to each other. While cosine similarity
is useful for clustering and pattern recognition, the projection approach
is generally more precise when it comes to quantifying actual contributions
and mixing between spin orbitals, especially in systems with open-shell
or near-degenerate character. As regards the diagonalization-based
projection strategy, it allows the detection of SOMO candidates with
negligible coupling to occupied $\beta$ orbitals, provides a detailed
inspection of how $\alpha$ occupied orbitals distribute onto the
$\beta$ manifold, and offers deeper insights  of orbital reorganization
effects in open-shell systems. 

\section{Examples}

\subsection{Formaldehyde (H$_{2}$CO) in its lowest triplet state}

In an all-electron basis set, there are 9 occupied $\alpha$ MOs,
$\phi_{\text{occ}}^{\alpha}$, and 7 occupied $\beta$ MOs, $\phi_{\text{occ}}^{\beta}$.
As summarized in the previous section, {\small\texttt{gSOMOs}} provides
several tools to find two two SOMOs among the nine $\phi_{\text{occ}}^{\alpha}$.
Table \ref{tab:Projection-of-alpha2beta-H2CO} presents simplified
projection data of occupied $\alpha$ orbitals onto $\beta$ orbitals
for the lowest triplet state, T\textsubscript{1}, of formaldehyde.
It is adapted from the dataframe created by the \texttt{project\_occupied\_alpha\_onto\_beta()}
function. Orbitals identified as SOMO indicate significant projection
onto virtual $\beta$ orbitals and negligible projection onto occupied
$\beta$ orbitals. 

\begin{sidewaystable}[H]
\centering{}{\small{}%
\begin{tabular}{ccccccccc}
\toprule 
{\small$\boldsymbol{\alpha}$}{\small\textbf{ MO }} & {\small\textbf{Occ }}{\small$\boldsymbol{\alpha}$}{\small\textbf{ }} & {\small\textbf{Energy (Ha) }} & {\small\textbf{P\texttwosuperior{} $\boldsymbol{\beta_{\mathrm{virt}}}$ }} & {\small\textbf{P\texttwosuperior{} $\boldsymbol{\beta_{\mathrm{occ}}}$ }} & {\small\textbf{$\boldsymbol{\beta}$ MO{*} }} & {\small\textbf{Occ $\boldsymbol{\beta}$ }} & {\small\textbf{SOMO? }} & \textbf{$\boldsymbol{\beta}$ MOs \textgreater 15\%}\tabularnewline
\midrule
\midrule 
\rowcolor{green!30}{\small 9} & {\small O } & {\small -0.204 } & {\small 0.996 } & {\small 0.004 } & {\small 9 } & {\small V } & {\small Y } & 9: 96.2\%\tabularnewline
\rowcolor{green!30}{\small 8} & {\small O } & {\small -0.367 } & {\small 0.895 } & {\small 0.105 } & {\small 8 } & {\small V } & {\small Y } & 8: 88.4\%\tabularnewline
{\small 7} & {\small O } & {\small -0.438 } & {\small 0.005 } & {\small 0.995 } & {\small 7 } & {\small O } & {\small N } & 7: 97.5\%\tabularnewline
{\small 6} & {\small O } & {\small -0.487 } & {\small 0.000 } & {\small 1.000 } & {\small 6 } & {\small O } & {\small N } & 6: 97.9\%\tabularnewline
{\small 5} & {\small O } & {\small -0.529 } & {\small 0.106 } & {\small 0.894 } & {\small 5 } & {\small O } & {\small N } & 5: 89.4\%\tabularnewline
{\small 4} & {\small O } & {\small -0.672 } & {\small 0.002 } & {\small 0.998 } & {\small 4 } & {\small O } & {\small N } & 4: 99.4\%\tabularnewline
{\small 3} & {\small O } & {\small -1.096 } & {\small 0.001 } & {\small 0.999 } & {\small 3 } & {\small O } & {\small N } & 3: 99.6\%\tabularnewline
{\small 2} & {\small O } & {\small -10.238 } & {\small 0.000 } & {\small 1.000 } & {\small 2 } & {\small O } & {\small N } & 2: 100.0\%\tabularnewline
{\small 1} & {\small O } & {\small -19.240 } & {\small 0.000 } & {\small 1.000 } & {\small 1 } & {\small O } & {\small N } & 1: 100.0\%\tabularnewline
\bottomrule
\end{tabular}}{\small\caption{Projection of $\alpha$ molecular orbitals onto $\beta$ space for
formaldehyde (H$_{2}$CO), highlighting SOMOs. {\small\textbf{P\texttwosuperior{}
$\boldsymbol{\beta_{\mathrm{virt}}}$}} and {\small\textbf{P\texttwosuperior{}
$\boldsymbol{\beta_{\mathrm{occ}}}$}} are $\|\mathbf{v}_{i}^{\text{occ}}\|^{2}$
and $\|\mathbf{v}_{i}^{\text{virt}}\|^{2}$ , respectively (see text).}\label{tab:Projection-of-alpha2beta-H2CO}
}
\end{sidewaystable}

A heatmap can ne generated to visualize the main projection contributions
of $\alpha$ molecular orbitals onto $\beta$ orbitals. The color
intensity reflects the percentage contribution of each $\beta$ orbital
to the total projection norm of a given $\alpha$ orbital. Only contributions
above 15\% were retained for clarity. Red dashed lines indicate the
HOMO--LUMO frontier for both spin channels (Figure \ref{fig:Similarity-of-alphaWithBeta-H2CO}a).
This analysis is very close to the heatmap generated after the cosine
similarity between $\alpha$ and $\beta$ MOs (Figure \ref{fig:Similarity-of-alphaWithBeta-H2CO}b)

\begin{figure}[h]
\begin{centering}
\includegraphics[width=8cm]{H2CO_T1_projection_heatmap-C}\includegraphics[width=8cm]{H2CO_T1_heatmap-C}
\par\end{centering}
\caption{similarity of $\alpha$ and $\beta$ MOs of the first triplet state
of H\protect\textsubscript{2}CO around the HOMO-LUMO frontier. (a)
Projection of $\alpha$ MOs onto the full space spanned by all $\beta$
orbitals (only contributions above 20\% were retained for clarity);
(b) Cosine similarity. Dashed lines mark the HOMO/LUMO boundaries
for $\alpha$ (horizontal) and $\beta$ (vertical) spin orbitals.}\label{fig:Similarity-of-alphaWithBeta-H2CO}

\end{figure}

And finally, and still in a projection scheme, the analysis of the
eigenvalues of $\boldsymbol{P}_{\text{occ}}$ and $\boldsymbol{P}_{\text{virt}}$
agree with the previous similarity analysis schemes. Figure \ref{fig:Eigenvalues-of-PoccPvirt-H2CO}
reveals that two occupied $\alpha$ MOs do not project on the $\phi_{\text{occ}}^{\beta}$
subspace, whereas two MOs - hopefully the same - project on the $\phi_{\text{virt}}^{\beta}$
subspace. The MO analysis performed on the two eigenvectors of $\boldsymbol{P}_{\text{occ}}$
and the two eigenvectors of $\boldsymbol{P}_{\text{virt}}$ and reported
in Table \ref{tab:SOMO-candidates-forH2CO} show without doubt that
the counterpart of the 8th and 9th $\alpha$ MOs are the 8th and 9th
unoccupied $\beta$ MOs.

\begin{figure}[h]
\begin{centering}
\includegraphics[width=16cm]{H2CO_T1_ProjectionEigenvalues-C}
\par\end{centering}
\caption{Eigenvalues of $\boldsymbol{P}_{\text{occ}}$ and $\boldsymbol{P}_{\text{virt}}$
for H\protect\textsubscript{2}CO.}\label{fig:Eigenvalues-of-PoccPvirt-H2CO}

\end{figure}

\begin{table}[h]
\begin{centering}
\begin{tabular}{cc}
{\small\texttt{\textbf{SOMO Candidate \#1}}} & {\small\texttt{\textbf{SOMO Candidate \#2}}}\tabularnewline
\hline 
\begin{cellvarwidth}[t]
\centering
{\small\texttt{$\alpha$ occupied contributions:}}{\small\par}

{\small\texttt{• $\alpha$ 9 (99.5\%)}}{\small\par}

{\small\texttt{$\beta$ virtual projections:}}{\small\par}

{\small\texttt{• $\beta$ 9 (96.4\%)}}
\end{cellvarwidth} & \begin{cellvarwidth}[t]
\centering
{\small\texttt{$\alpha$ occupied contributions:}}{\small\par}

{\small\texttt{• $\alpha$ 8 (89.4\%)}}{\small\par}

{\small\texttt{$\beta$ virtual projections:}}{\small\par}

{\small\texttt{• $\beta$ 8 (98.4\%)}}
\end{cellvarwidth}\tabularnewline
\hline 
\end{tabular}
\par\end{centering}
\caption{SOMO candidates for H\protect\textsubscript{2}CO, as given by the
analyzis of the eigenvectors of $\boldsymbol{P}_{\text{occ}}$ and
$\boldsymbol{P}_{\text{virt}}$.}\label{tab:SOMO-candidates-forH2CO}

\end{table}

\pagebreak{}

\subsection{Iron complex}

\begin{wrapfigure}{O}{0.5\columnwidth}%
\begin{centering}
\includegraphics[width=7cm]{FeComplex-C}
\par\end{centering}
\caption{Iron complex in its fully optimized quintet state geometry.}\label{fig:Iron-complex-geom}
\end{wrapfigure}%
In some cases, the identification of SOMOs by visual inspection can
become very challenging, making a tool like {\small\texttt{gSOMOs}}
particularly valuable. The dispersion of SOMO contributions over multiple
$\alpha$ and $\beta$ orbitals can arise from near-degeneracy effects
in the frontier orbital region, possibly amplified by minor symmetry
breaking or delocalization effects inherent to DFT. This behavior
is typical for open-shell transition metal complexes with dense manifolds
of occupied and virtual states This is the case for the quintet state
of the iron complex shown in Figure \ref{fig:Iron-complex-geom}.
The cosine similarity method correctly identifies three SOMOs but
fails for the fourth, which only imperfectly projects onto the virtual
$\beta$ space - whereas it correctly identifies MOs 169, 92 and 194,
the fourth SOMOs is identified as MO 164. In contrast, the projection
scheme performs well, especially through the diagonalization of $\boldsymbol{P}_{\text{occ}}$
and $\boldsymbol{P}_{\text{virt}}$\LyXZeroWidthSpace{} and the associated
analysis. Figure \ref{fig:Eigenvalues-of-PoccPvirt-FeComplex} shows
that four null eigenvalues are found after the diagonalization of
$\boldsymbol{P}_{\text{occ}}$\LyXZeroWidthSpace , associated to four
eigenvalues close to 1 after the diagonalization of $\boldsymbol{P}_{\text{virt}}$.
The decomposition of the SOMO candidates is reported in Table \ref{tab:SOMO-candidates-for-FeComplex}.
The projection scheme isolates clear contributions for each SOMO candidate,
mainly involving a few $\alpha$-occupied and $\beta$-virtual orbitals.
SOMO\textsubscript{1} shows a mixed $\alpha$-character (orbitals
187 and 164) projecting mainly onto $\beta$-orbital 194. SOMO\textsubscript{2}
has a more distributed $\alpha$-character but projects onto $\beta$-orbitals
192 and 193. SOMO\textsubscript{3} is dominated by orbital 186, projecting
onto $\beta$-orbital 198, while SOMO\textsubscript{4} is mainly
from orbital 168, with projections onto $\beta$-orbitals 193 and
192. SOMO\textsubscript{1} clearly shows why the cosine similarity
misidentified the dominant $\alpha_{\mathrm{occ}}$-$\beta_{\mathrm{virt}}$
pair: although orbital 187 contributes most (44.2\%), a significant
mixing with orbital 164 (27.3\%) leads to an overemphasis on the 164
→ 194 projection. This explains why the cosine similarity method failed
here, while the projection-based approach provides a more reliable
identification. Despite some mixing, the dominant contributions are
clearly identified, confirming the robustness of the projection analysis. 

\begin{figure}[h]
\begin{centering}
\includegraphics[width=16cm]{FeComplex_ProjectionEigenvalues-C}
\par\end{centering}
\caption{Eigenvalues of $\boldsymbol{P}_{\text{occ}}$ and $\boldsymbol{P}_{\text{virt}}$
for the iron complex.}\label{fig:Eigenvalues-of-PoccPvirt-FeComplex}
\end{figure}

\begin{table}[h]
\begin{centering}
\begin{tabular}{cc}
{\small\texttt{\textbf{SOMO Candidate \#1}}} & {\small\texttt{\textbf{SOMO Candidate \#2}}}\tabularnewline
\hline 
\begin{cellvarwidth}[t]
\centering
{\small\texttt{$\alpha$ occupied contributions:}}{\small\par}

{\small\texttt{• $\alpha$ 187 (44.2\%)}}{\small\par}

{\small\texttt{• $\alpha$ 164 (27.3\%)}}{\small\par}

{\small\texttt{$\beta$ virtual projections:}}{\small\par}

{\small\texttt{• $\beta$ 194 (73.3\%)}}{\small\par}

{\small\texttt{• $\beta$ 196 (16.1\%)}}
\end{cellvarwidth} & \begin{cellvarwidth}[t]
\centering
{\small\texttt{$\alpha$ occupied contributions:}}{\small\par}

{\small\texttt{• $\alpha$ 169 (41.1\%)}}{\small\par}

{\small\texttt{• $\alpha$ 186 (21.6\%)}}{\small\par}

{\small\texttt{• $\alpha$ 165 (15.7\%)}}{\small\par}

{\small\texttt{$\beta$ virtual projections:}}{\small\par}

{\small\texttt{• $\beta$ 192 (53.1\%)}}{\small\par}

{\small\texttt{• $\beta$ 193 (26.9\%)}}
\end{cellvarwidth}\tabularnewline
\hline 
\hline 
{\small\texttt{\textbf{SOMO Candidate \#3}}} & {\small\texttt{\textbf{SOMO Candidate \#4}}}\tabularnewline
\hline 
\begin{cellvarwidth}[t]
\centering
{\small\texttt{$\alpha$ occupied contributions:}}{\small\par}

{\small\texttt{• $\alpha$ 186 (30.0\%)}}{\small\par}

{\small\texttt{$\beta$ virtual projections:}}{\small\par}

{\small\texttt{• $\beta$ 198 (73.0\%)}}
\end{cellvarwidth} & \begin{cellvarwidth}[t]
\centering
{\small\texttt{$\alpha$ occupied contributions:}}{\small\par}

{\small\texttt{• $\alpha$ 168 (51.8\%)}}{\small\par}

{\small\texttt{• $\alpha$ 183 (16.3\%)}}{\small\par}

{\small\texttt{$\beta$ virtual projections:}}{\small\par}

{\small\texttt{• $\beta$ 193 (41.6\%)}}{\small\par}

{\small\texttt{• $\beta$ 192 (26.7\%)}}
\end{cellvarwidth}\tabularnewline
\hline 
\end{tabular}
\par\end{centering}
\caption{SOMO candidates for the ircon complex, as given by the analyzis of
the eigenvectors of $\boldsymbol{P}_{\text{occ}}$ and $\boldsymbol{P}_{\text{virt}}$.}\label{tab:SOMO-candidates-for-FeComplex}
\end{table}


\section{Conclusion}
\end{document}
